% =====================================================================
%  BÖLÜM 11: İstatistiğin Temelleri — Model ve Örneklem
%  Kaynak: Casella & Berger; Hogg & Craig; Lehmann & Casella
% =====================================================================
\chapter{İstatistiğin Temelleri: Model ve Örneklem}
\label{chp:istatistik_temelleri}

\bolumalinti{%
  İstatistik, gözlemlerin ötesinde olanı anlamaya çalışır;
  bu anlamda her istatistiksel çıkarım bir inançtır.%
}{Erich Leo Lehmann, \textit{Theory of Point Estimation}}

% ---- Kazanımlar (TYMM) ----
\begin{kazanimlar}
Bu bölümün sonunda öğrenci:
\begin{itemize}[leftmargin=1.8em]
  \item İstatistiksel model kavramını parametrik ve parametrik olmayan çerçevede tanımlayabilir;
        istatistiksel karar, kayıp fonksiyonu ve risk kavramlarını açıklayabilir.
  \item Örneklem, istatistik ve örneklem dağılımı kavramlarını tanımlayabilir;
        $\bar{X}_n$ ve $S^2$'nin dağılımlarını türetebilir.
  \item Yeterli istatistiği tanımlayabilir; Fisher-Neyman çarpanlaştırma teoremini
        ispatlayabilir ve uygulamalarda yeterli istatistik bulabilir.
  \item Tamamlayıcı ve en az yeterli istatistik kavramlarını tanımlayabilir;
        bunları UMVUE teorisindeki rolleriyle ilişkilendirebilir.
\end{itemize}
\end{kazanimlar}

% ---- Motivasyon ----
\begin{motivasyon}
Part~I'de olasılığın matematiksel temellerini inşa ettik: ölçüm teorisi,
rasgele değişkenler, beklenti, dağılımlar, büyük sayılar ve limit teoremler.
Şimdi bu araçları ters yönde kullanacağız: \emph{gözlemlerden} hareketle
bilinmeyen parametreler hakkında çıkarım yapmak.

İstatistik dili şöyle çalışır: gerçeklik, bilinmeyen parametre $\theta$ tarafından
belirlenen bir dağılım $P_\theta$'dur. Biz yalnızca bu dağılımdan çekilen
$n$ gözlemi görürüz ve $\theta$'yı tahmin etmeye ya da test etmeye çalışırız.

\medskip
\textbf{Köprü.} Bölüm~\ref{chp:dagilim_aileleri}'deki üstel aile,
Bölüm~\ref{chp:beklenti}'deki koşullu beklenti ve
Bölüm~\ref{chp:merkezi_limit}'teki asimptotik teori bu bölümün doğrudan ön koşullarıdır.
\end{motivasyon}

% =====================================================================
\section{İstatistiksel Model Kavramı}
\label{sec:istatistik_model}
% =====================================================================

\begin{tanim}[İstatistiksel Model]
\label{def:istatistik_model}
\textbf{İstatistiksel model} bir aile $\mathcal{P} = \{P_\theta : \theta \in \Theta\}$'dır;
burada her $P_\theta$ $(\mathcal{X}, \mathcal{B})$ üzerinde bir olasılık ölçüsü,
$\Theta \subseteq \mathbb{R}^k$ parametre uzayıdır.

\begin{itemize}
  \item \textbf{Parametrik model:} $\Theta$ sonlu boyutlu ($\Theta \subseteq \mathbb{R}^k$, $k$ sabit).
        Örnek: $\Normal(\mu, \sigma^2)$ ailesi, $\Theta = \mathbb{R}\times\mathbb{R}_{>0}$.
  \item \textbf{Parametrik olmayan model:} $\Theta$ sonsuz boyutlu.
        Örnek: "tüm sürekli dağılımlar" ailesi.
  \item \textbf{Yarı-parametrik model:} Parametrik bileşen + parametrik olmayan bileşen.
        Örnek: Cox orantılı tehlike modeli.
\end{itemize}
\end{tanim}

\begin{tanim}[Tanımlanabilirlik]
\label{def:tanimlanabilirlik}
Model $\{P_\theta\}$ \textbf{tanımlanabilir} (identifiable) dır, eğer
$\theta \neq \theta'$ $\Rightarrow$ $P_\theta \neq P_{\theta'}$ ise.
Tanımlanabilirlik olmadan $\theta$'yı veri ile ayırt etmek imkânsızdır.
\end{tanim}

\begin{ornek}[Tanımlanabilir ve Tanımlanamaz Modeller]
\textbf{Tanımlanabilir:} $\Normal(\mu, \sigma^2)$: farklı $(\mu,\sigma^2)$ çiftleri
farklı dağılımlar üretir.\\
\textbf{Tanımlanamaz (örnek):} $\Normal(\mu_1+\mu_2, \sigma^2)$:
$\mu_1=1, \mu_2=2$ ile $\mu_1=0, \mu_2=3$ aynı dağılımı verir.
Bu model $\mu_1$ ve $\mu_2$'yi ayrı ayrı tanımlayamaz.
\end{ornek}

% =====================================================================
\section{İstatistiksel Kararlar ve Kayıp Fonksiyonu}
\label{sec:karar_teorisi}
% =====================================================================

\begin{tanim}[Karar Çerçevesi]
\label{def:karar}
\textbf{Karar teorisi} çerçevesi:
\begin{itemize}
  \item $\Theta$: parametre uzayı (gerçeklik kümesi),
  \item $\mathcal{D}$: karar uzayı (verilen yanıtlar kümesi),
  \item $L(\theta, d)$: \textbf{kayıp fonksiyonu} (loss function) —
        gerçek parametre $\theta$ iken $d$ kararı almanın maliyeti,
  \item $\delta(\mathbf{X})$: \textbf{karar kuralı} — gözlemden karara giden fonksiyon.
\end{itemize}
\textbf{Risk fonksiyonu:} $R(\theta, \delta) = E_\theta[L(\theta, \delta(\mathbf{X}))]$.
\end{tanim}

\begin{ornek}[Yaygın Kayıp Fonksiyonları]
\begin{itemize}
  \item \textbf{Karesel kayıp (squared error loss):} $L(\theta, d) = (\theta - d)^2$.
        Risk $= E[(\theta - \delta(\mathbf{X}))^2]$ = ortalama karesel hata (OKH, MSE).
  \item \textbf{Mutlak değer kaybı:} $L(\theta, d) = |\theta - d|$.
        Optimal karar: koşullu medyan.
  \item \textbf{0-1 kaybı (sınıflandırma):} $L(\theta, d) = \mathbf{1}_{d\neq\theta}$.
        Risk = hata olasılığı.
\end{itemize}
\end{ornek}

\begin{tanim}[Yansızlık ve OKH Ayrışımı]
\label{def:yansizlik_okh}
$\delta(\mathbf{X})$, $\theta$'nın \textbf{yansız tahmin edicisi} (unbiased estimator):
$E_\theta[\delta(\mathbf{X})] = \theta$ her $\theta \in \Theta$.
Karesel kayıp altında:
\[
  \mathrm{OKH}(\delta, \theta) = \mathrm{Var}_\theta(\delta) + [\mathrm{Yan}(\delta, \theta)]^2,
\]
burada $\mathrm{Yan}(\delta,\theta) = E_\theta[\delta] - \theta$. Yansız: $\mathrm{OKH} = \mathrm{Var}$.
\end{tanim}

\begin{proof}[OKH Ayrışımı]
$E[(\delta-\theta)^2] = E[(\delta - E[\delta] + E[\delta]-\theta)^2]
= \mathrm{Var}(\delta) + (E[\delta]-\theta)^2$ (çapraz terim sıfır).
\end{proof}

% =====================================================================
\section{Örneklem ve İstatistik Kavramı}
\label{sec:orneklem}
% =====================================================================

\begin{tanim}[Rasgele Örneklem]
\label{def:rasgele_orneklem}
$X_1, \ldots, X_n$ i.i.d.\ $P_\theta$ dağılımlı rasgele değişkenler;
bu dizi $\theta$'lı dağılımdan alınan \textbf{rasgele örneklem} (random sample) dir.
Veri $\mathbf{x} = (x_1,\ldots,x_n)$ bu dizinin bir gerçekleşmesidir.
\end{tanim}

\begin{tanim}[İstatistik]
\label{def:istatistik}
$T = T(X_1,\ldots,X_n) : \mathcal{X}^n \to \mathbb{R}^k$ ölçülebilir bir fonksiyon;
$\theta$'ya \emph{doğrudan} bağlı olmayan (parametreyi bilmeden hesaplanabilen) böyle
bir fonksiyona \textbf{istatistik} (statistic) denir.

Temel istatistikler:
\begin{align*}
  \bar{X}_n &= \frac{1}{n}\sum_{i=1}^n X_i \quad \text{(örneklem ortalaması)}, \\
  S^2 &= \frac{1}{n-1}\sum_{i=1}^n (X_i - \bar{X}_n)^2 \quad \text{(örneklem varyansı)}, \\
  M_k &= \frac{1}{n}\sum_{i=1}^n X_i^k \quad \text{($k$'ıncı örneklem momenti)}.
\end{align*}
\end{tanim}

% =====================================================================
\section{Örneklem Dağılımları}
\label{sec:orneklem_dagilim}
% =====================================================================

\subsection{Örneklem Ortalaması ve Varyansının Dağılımı}

\begin{teorem}[Normal Popülasyondan Örneklem]
\label{thm:normal_orneklem}
$X_1,\ldots,X_n$ i.i.d.\ $\Normal(\mu,\sigma^2)$. O zaman:
\begin{enumerate}[(i)]
  \item $\bar{X}_n \sim \Normal(\mu,\, \sigma^2/n)$.
  \item $\frac{(n-1)S^2}{\sigma^2} \sim \chi^2_{n-1}$.
  \item $\bar{X}_n$ ve $S^2$ \textbf{bağımsızdır}.
  \item $T = \frac{\bar{X}_n - \mu}{S/\sqrt{n}} \sim t_{n-1}$.
\end{enumerate}
\end{teorem}

\begin{proof}
\textit{(i):} $\bar{X}_n = n^{-1}\sum X_i$; bağımsız normalların doğrusal bileşimi
Bölüm~\ref{thm:normal_yeniden}'den $\Normal(\mu, \sigma^2/n)$.

\textit{(ii) ve (iii):} Ortogonal dönüşüm argümanı.
$(X_1,\ldots,X_n)$ birim vektörlerle ifade edilirse, $\bar{X}_n$
$\mathbf{1}/\sqrt{n}$ yönüne, $\sum(X_i-\bar{X}_n)^2$ ise buna dik olan $n-1$ boyutlu
alt uzaya karşılık gelir. Ortogonal dönüşüm normalliği korur ve bağımsızlığı üretir.
Dik bileşen: $n-1$ bağımsız standart normal karesinin toplamı $= \chi^2_{n-1}$.

Alternatif ispat: $Y_1 = \sqrt{n}\bar{X}$, $Y_i = X_{i-1}-X_i$ ($i=2,\ldots,n$) dönüşümü;
Jacobian hesabı; Fubini ile bağımsızlık.

\textit{(iv):} $Z = \sqrt{n}(\bar{X}_n-\mu)/\sigma \sim \Normal(0,1)$;
$V = (n-1)S^2/\sigma^2 \sim \chi^2_{n-1}$ bağımsız.
$T = Z/\sqrt{V/(n-1)} \sim t_{n-1}$ (tanım gereği).
\end{proof}

\begin{hataanalizi}
"$S^2 = \frac{1}{n}\sum(X_i-\bar{X}_n)^2$" örneklem varyansı için yaygın bir notasyon.
Ancak bu yanlı tahmin edicidir: $E[S_n^2] = (n-1)\sigma^2/n \neq \sigma^2$.
Kitabımızda $S^2 = \frac{1}{n-1}\sum(X_i-\bar{X}_n)^2$ (Bessel düzeltmesi) kullanılır:
$E[S^2] = \sigma^2$.
\end{hataanalizi}

\begin{teorem}[Örneklem Ortalamasının Büyük Örneklem Dağılımı]
\label{thm:buyuk_orneklem_dagilim}
$X_1,\ldots,X_n$ i.i.d., $E[X_i]=\mu$, $\mathrm{Var}(X_i)=\sigma^2<\infty$.
Dağılımı ne olursa olsun: $\bar{X}_n \xrightarrow{d} \Normal(\mu, \sigma^2/n)$ (MLT).
$S^2 \xrightarrow{P} \sigma^2$ (BSK + sürekli haritalama).
$T_n = \sqrt{n}(\bar{X}_n-\mu)/S \xrightarrow{d} \Normal(0,1)$ (Slutsky).
\end{teorem}

% =====================================================================
\section{Yeterli İstatistik}
\label{sec:yeterli_istatistik}
% =====================================================================

\begin{kesif}
$n$ gözlemden elde edilen tüm bilgi için $T(X_1,\ldots,X_n)$'yi hesaplamak yeterli mi?
Yani $T$'yi bilince, orijinal veriyi $\theta$ hakkında daha fazla bilgi için kullanabilir miyiz?
Yanıt "hayır" ise $T$ yeterli istatistik adını alır.
\end{kesif}

\begin{tanimgerekce}
$X_1,\ldots,X_n$ birlikte $\theta$ hakkında tüm bilgiyi içerir.
$T$ "yeterli"yse, $T$'nin değeri bilindiğinde verinin $\theta$'dan bağımsız bir kalan dağılımı olur.
Başka deyişle $T$, $\theta$ bilgisini "özetlemiş"tir ve orijinal veriden geriye kalan
rasgelelik $\theta$ ile ilgisizdir.
\end{tanimgerekce}

\begin{tanim}[Yeterli İstatistik]
\label{def:yeterli}
$T = T(\mathbf{X})$ istatistiği $\theta$ için \textbf{yeterli} (sufficient) dır, eğer
$\mathbf{X}$'in koşullu dağılımı $T = t$ verildiğinde $\theta$'dan bağımsız ise, yani
$P_\theta(\mathbf{X} \in A \mid T = t)$ her $t$ için $\theta$'ya bağlı değilse.
\end{tanim}

\begin{teorem}[Fisher-Neyman Çarpanlaştırma Teoremi]
\label{thm:fisher_neyman}
$\mathbf{X}$ ortak PDF/PMF'si $f(\mathbf{x} \mid \theta)$ olan rasgele vektör.
$T(\mathbf{X})$, $\theta$ için yeterlidir ancak ve ancak
\[
  f(\mathbf{x} \mid \theta) = g(T(\mathbf{x}), \theta) \cdot h(\mathbf{x})
\]
ayrışımı mevcut ise; burada $g$ yalnızca $T(\mathbf{x})$ ve $\theta$'ya, $h$ yalnızca $\mathbf{x}$'e bağlıdır.
\end{teorem}

\begin{proof}
\textbf{Kesikli durum ($\Leftarrow$):}
$P_\theta(\mathbf{X}=\mathbf{x} \mid T(\mathbf{X})=t) = \frac{P_\theta(\mathbf{X}=\mathbf{x},\, T(\mathbf{X})=t)}{P_\theta(T(\mathbf{X})=t)}$.
$T(\mathbf{x})=t$ ise pay $= f(\mathbf{x}|\theta) = g(t,\theta)h(\mathbf{x})$.
Payda $= \sum_{\mathbf{x}': T(\mathbf{x}')=t} g(t,\theta)h(\mathbf{x}') = g(t,\theta)\sum_{\mathbf{x}':T=t}h(\mathbf{x}')$.
Oran $= h(\mathbf{x})/\sum_{\mathbf{x}':T=t}h(\mathbf{x}')$ — $\theta$'dan bağımsız. $\Rightarrow$ $T$ yeterli.

\textbf{($\Rightarrow$):} $T$ yeterli olsun.
$f(\mathbf{x}|\theta) = P_\theta(\mathbf{X}=\mathbf{x}, T=T(\mathbf{x}))
= P_\theta(\mathbf{X}=\mathbf{x}|T=T(\mathbf{x})) \cdot P_\theta(T=T(\mathbf{x}))$.
$P_\theta(\mathbf{X}=\mathbf{x}|T=T(\mathbf{x})) =: h(\mathbf{x})$ ($\theta$'dan bağımsız).
$P_\theta(T=T(\mathbf{x})) =: g(T(\mathbf{x}),\theta)$.
Ayrışım sağlandı.
\end{proof}

\begin{ornek}[Binom için Yeterli İstatistik]
$X_1,\ldots,X_n$ i.i.d.\ $\mathrm{Bernoulli}(p)$.
$f(\mathbf{x}|p) = \prod p^{x_i}(1-p)^{1-x_i} = p^{\sum x_i}(1-p)^{n-\sum x_i}$.
$g(T,p) = p^T(1-p)^{n-T}$ ($T = \sum x_i$), $h(\mathbf{x}) = 1$.
Fisher-Neyman: $T = \sum X_i$ yeterli.
\end{ornek}

\begin{ornek}[Normal için Yeterli İstatistik]
$X_i \sim \Normal(\mu, \sigma^2)$, $\theta = (\mu, \sigma^2)$.
$f(\mathbf{x}|\theta) = (2\pi\sigma^2)^{-n/2}\exp\!\left(-\frac{\sum(x_i-\mu)^2}{2\sigma^2}\right)$.
$\sum(x_i-\mu)^2 = \sum x_i^2 - 2\mu\sum x_i + n\mu^2$.
$g(T_1,T_2,\theta) = (2\pi\sigma^2)^{-n/2}\exp\!\left(\frac{2\mu T_1 - T_2 - n\mu^2}{2\sigma^2}\right)$,
$h(\mathbf{x})=1$; $T_1 = \sum x_i$, $T_2 = \sum x_i^2$.
$(\sum X_i, \sum X_i^2)$ — dolayısıyla $(\bar{X}_n, S^2)$ — yeterli.
\end{ornek}

\begin{ornek}[Üstel Dağılım]
$X_i \sim \mathrm{Exp}(\lambda)$: $f(\mathbf{x}|\lambda) = \lambda^n e^{-\lambda\sum x_i}$.
$T = \sum X_i$ yeterli; $g(T,\lambda) = \lambda^n e^{-\lambda T}$, $h = 1$.
\end{ornek}

\begin{teorem}[Üstel Ailede Yeterli İstatistik]
\label{thm:ustel_yeterli}
Tek parametreli üstel aile $f(x|\theta) = h(x)\exp(\eta(\theta)T(x) - B(\theta))$.
$n$ gözlemde: $f(\mathbf{x}|\theta) = \prod h(x_i) \cdot \exp\!\bigl(\eta(\theta)\sum T(x_i) - nB(\theta)\bigr)$.
Fisher-Neyman: $\sum_{i=1}^n T(X_i)$ yeterli istatistiktir.
\end{teorem}

\begin{mikroozet}
Yeterlilik $\equiv$ çarpanlaştırma: $f = g(T,\theta)\cdot h(\mathbf{x})$.
Üstel ailede her zaman doğal yeterli istatistik: $\sum T(X_i)$.
\end{mikroozet}

% =====================================================================
\section{Tamamlayıcı İstatistik}
\label{sec:tamamlayici}
% =====================================================================

\begin{tanim}[Tamamlayıcı İstatistik]
\label{def:tamamlayici}
$T$ istatistiği $\theta$ için \textbf{tamamlayıcı} (complete) dır, eğer
her ölçülebilir $g : \mathbb{R} \to \mathbb{R}$ için
\[
  E_\theta[g(T)] = 0 \quad \forall\,\theta \in \Theta \implies g(T) = 0 \quad \text{n.k.\ (her } \theta\text{'da)}.
\]
\end{tanim}

\begin{onerme}[Tamamlayıcılığın Önemi]
Tamamlayıcı bir istatistik varsa, $T$ üzerinden ifade edilen her yansız tahmin edici tek türlüdür.
Bu, UMVUE (en küçük varyanslı yansız tahmin edici) teorisinin kilit adımıdır
(Bölüm~\ref{chp:nokta_tahmin}, Lehmann-Scheffé teoremi).
\end{onerme}

\begin{teorem}[Üstel Ailede Tamamlayıcılık]
\label{thm:ustel_tam}
Tam rankli (tam parametreli) tek parametreli üstel ailede yeterli istatistik
$T = \sum T(X_i)$ tamamlayıcıdır.
\end{teorem}

\begin{proof}
$E_\theta[g(T)] = 0$ tüm $\theta$'da koşulu, Laplace ya da Moment dönüşümü argümanıyla
$g = 0$ a.e.\ verir. Üstel aile ailesiyle MÜF biricikliği kullanılır.
Ayrıntı: Lehmann-Casella~\cite{Lehmann1998}, Teorem 1.6.22.
\end{proof}

\begin{ornek}[Binom Tamamlayıcılığı]
$T = \sum X_i \sim \Binom(n,p)$. $E_p[g(T)] = 0$ tüm $p \in (0,1)$:
$\sum_{k=0}^n g(k)\binom{n}{k}p^k(1-p)^{n-k} = 0$.
Bu $p$ polinomu — sıfır polinom ancak $g(k) = 0$ tüm $k$'larda.
$T$ tamamlayıcı.
\end{ornek}

% =====================================================================
\section{En Az Yeterli İstatistik}
\label{sec:en_az_yeterli}
% =====================================================================

\begin{tanim}[En Az Yeterli İstatistik]
\label{def:en_az_yeterli}
$T$ \textbf{en az yeterli} (minimal sufficient) dır, eğer:
\begin{enumerate}[(i)]
  \item $T$ yeterlidir.
  \item Her yeterli $U$ için, $U$'nun bir fonksiyonu olarak $T$ ifade edilebilir:
        $T = h(U)$ (ölçülebilir $h$).
\end{enumerate}
En az yeterli istatistik, veriyi mümkün olan en yoğun biçimde özetler.
\end{tanim}

\begin{teorem}[Lehmann-Scheffé En Az Yeterlilik Kriteri]
\label{thm:ls_en_az}
$f(\mathbf{x}|\theta)/f(\mathbf{y}|\theta)$ oranının $\theta$'dan bağımsız olması
$\mathbf{x}$ ve $\mathbf{y}$'nin aynı $T$ değerini vermesiyle eşdeğer ise,
$T$ en az yeterlidir.
\end{teorem}

\begin{ornek}[Normal — En Az Yeterlilik]
$X_i \sim \Normal(\mu, \sigma^2)$.
$f(\mathbf{x}|\theta)/f(\mathbf{y}|\theta) = \exp\!\bigl(\frac{\sum x_i^2-\sum y_i^2}{-2\sigma^2}
+ \frac{\mu(\sum x_i - \sum y_i)}{\sigma^2}\bigr)$.
Bu $\theta$'dan bağımsız $\Leftrightarrow$ $\sum x_i = \sum y_i$ ve $\sum x_i^2 = \sum y_i^2$.
$(\sum X_i, \sum X_i^2)$ en az yeterli (= tamamlayıcı yeterli $\Rightarrow$ UMVUE mevcut).
\end{ornek}

\begin{hocanotu}
Tamamlayıcı yeterli istatistik (complete sufficient statistic) her zaman en az yeterlidir.
Tersine, en az yeterli istatistik her zaman tamamlayıcı olmak zorunda değildir —
ama üstel aileler için ikisi çakışır.
\end{hocanotu}

% ---- Disiplinlerarası Bağlantı (TYMM) ----
\begin{kopru}[Disiplinlerarası — Veri Bilimi, Biyomedikal Araştırma, Ekonometri]
\textbf{Veri bilimi:} "Hangi özellikler (features) modeli tam olarak özetler?" sorusu,
yeterli istatistik kavramının makine öğrenmesi karşılığıdır.
Sıkıştırma (compression) teorisi, yeterlilik ile entropi arasındaki bağı kurar.\\
\textbf{Klinik araştırma:} İlaç klinik denemelerinde birincil uç nokta (primary endpoint)
seçimi bir istatistik seçim problemidir; yeterli ve duyarlı bir istatistik
çalışmanın güç (power) hesabını doğrudan etkiler.\\
\textbf{Ekonometri:} Araç değişken (instrumental variable) ve panel veri modellerinde
parametre tanımlanabilirliği, bu bölümdeki kavramlarla doğrudan ilişkilidir.
\end{kopru}

% ---- Öz-Yansıtma ----
\begin{ozyansima}
\begin{itemize}[leftmargin=1.8em]
  \item "Parametrik model" ve "parametrik olmayan model" arasındaki farkı bir örnekle açıklayın.
  \item Yeterli istatistiğin "veri özetleme" yorumunu ve pratikte neden önemli olduğunu anlatın.
  \item $(\bar{X}_n, S^2)$'nin $\sigma^2$ bilinen $\Normal(\mu,\sigma^2)$'de yeterli istatistik olup
        olmadığını tartışın. Tek bilinmeyen $\mu$ ise ne değişir?
  \item OKH ayrışımı $= \mathrm{Var} + \mathrm{Yan}^2$ istatistiksel bir "uzlaşı"yı nasıl yansıtır?
\end{itemize}
\end{ozyansima}

% ---- Köprü Notu ----
\begin{kopru}
\textbf{Önceki kısım:} Part~I'deki dağılım teorisi (Böl.~\ref{chp:dagilim_aileleri}),
MLT (Böl.~\ref{chp:merkezi_limit}) ve koşullu beklenti (Böl.~\ref{chp:beklenti})
bu bölümde doğrudan kullanıldı.\\
\textbf{Sonraki bölüm:} Bölüm~\ref{chp:nokta_tahmin} — yeterli istatistik + tamamlayıcılık
$\to$ UMVUE; Fisher bilgisi + Cramér-Rao $\to$ verimlilik alt sınırı; MLE.\\
\textbf{Bölüm~\ref{chp:guven_araliklari}:} Pivot yöntemi örneklem dağılımlarına dayanır.
\defterbaglantisi{bolum11\_istatistik\_temelleri.ipynb}{%
  Örneklem dağılımı simülasyonu, yeterli istatistik görselleştirmesi}
\end{kopru}

% =====================================================================
\section{Sıra İstatistikleri}
\label{sec:sira_istatistik}
% =====================================================================

\begin{tanim}[Sıra istatistikleri]\label{def:sira_ist}
$X_1,\ldots,X_n\iid F$ sürekli dağılım. Sıralı değerler $X_{(1)}\leq X_{(2)}\leq\cdots\leq X_{(n)}$ \textbf{sıra istatistikleri} \emph{(order statistics)}, toplu vektörü $(X_{(1)},\ldots,X_{(n)})$ \textbf{sıralanmış örneklem} adını alır.
\end{tanim}

\begin{teorem}[Sıra istatistiği yoğunluğu]\label{thm:sira_yogunluk}
$X_{(k)}$'nın yoğunluğu:
\[
  f_{X_{(k)}}(x) = \frac{n!}{(k-1)!(n-k)!}[F(x)]^{k-1}[1-F(x)]^{n-k}f(x).
\]
$X_{(1)}$ için $k=1$: $f_{X_{(1)}}(x) = n[1-F(x)]^{n-1}f(x)$.
$X_{(n)}$ için $k=n$: $f_{X_{(n)}}(x) = n[F(x)]^{n-1}f(x)$.
\end{teorem}

\begin{proof}
Kombinatorik argüman: $P(X_{(k)}\in[x,x+dx])$ = $\binom{n}{1}\binom{n-1}{k-1}$ yolu ile $k-1$ gözlem $<x$, bir gözlem $\in[x,x+dx]$, $n-k$ gözlem $>x$.
\end{proof}

\begin{teorem}[Min-max için özel formüller]\label{thm:min_max}
$X_i\iid F$ bağımsız.
\begin{enumerate}[(a)]
  \item $F_{X_{(1)}}(x) = 1-[1-F(x)]^n$ (minimum KDF).
  \item $F_{X_{(n)}}(x) = [F(x)]^n$ (maksimum KDF).
  \item Aralık (range): $R_n=X_{(n)}-X_{(1)}$.
\end{enumerate}
\end{teorem}

\begin{ornek}[Uniform sıra istatistikleri]\label{ex:uniform_sira}
$X_i\iid\Uniform[0,1]$. $X_{(k)}\sim\mathrm{Beta}(k,n-k+1)$; $E[X_{(k)}]=k/(n+1)$. Dolayısıyla sıra istatistiklerinin beklentileri $[0,1]$'i eşit aralıklara böler.
\end{ornek}

% =====================================================================
\section{Ampirik Dağılım Fonksiyonu}
\label{sec:ampirik_dagilim}
% =====================================================================

\begin{tanim}[Ampirik dağılım fonksiyonu]\label{def:edf}
$X_1,\ldots,X_n\iid F$. \textbf{Ampirik dağılım fonksiyonu} \emph{(empirical distribution function, EDF)}:
\[
  F_n(x) = \frac{1}{n}\sum_{i=1}^n\mathbf{1}[X_i\leq x].
\]
$F_n$ sıra istatistiklerinin basamak fonksiyonu: $F_n(X_{(k)}) = k/n$.
\end{tanim}

\begin{teorem}[Glivenko-Cantelli]\label{thm:glivenko_cantelli}
$\|F_n - F\|_\infty = \sup_x|F_n(x)-F(x)|\asto0$.
\end{teorem}

\begin{proof}[Proof (taslak)]
Her sabit $x$ için $F_n(x)\asto F(x)$ (GBSK). Süreklilik + sonlu sayıda bölüm argümanıyla süpremum kontrolü sağlanır. Detay: Bölüm~\ref{chp:buyuk_sayilar}.
\end{proof}

\begin{teorem}[Kolmogorov-Smirnov]\label{thm:ks_istatistik}
$F$ sürekli altında $\sqrt{n}\|F_n-F\|_\infty\dto\|\mathbb{B}\|_\infty$ (Bölüm~\ref{chp:asimptotik}, Donsker teoremi). $H_0:F=F_0$ testi:
\[
  D_n = \sqrt{n}\sup_x|F_n(x)-F_0(x)|.
\]
$\alpha=0.05$ kritik değeri $\approx1.36$.
\end{teorem}

% =====================================================================
\section{Neyman-Fisher Çarpanlaştırma — Daha Fazla Örnek}
\label{sec:carpanlastirma_ornekler}
% =====================================================================

\begin{ornek}[İki parametreli Normal — $\mu$ bilinmiyor, $\sigma^2$ sabit]\label{ex:normal_mu_yeterli}
$X_i\iid\Normal(\mu,\sigma_0^2)$. Fisher-Neyman:
\[
  f(\mathbf{x}|\mu) = (2\pi\sigma_0^2)^{-n/2}\exp\!\left(-\frac{\sum x_i^2}{2\sigma_0^2}+\frac{\mu\sum x_i}{\sigma_0^2}-\frac{n\mu^2}{2\sigma_0^2}\right).
\]
$g(T,\mu)=\exp(\mu T/\sigma_0^2-n\mu^2/(2\sigma_0^2))$ ($T=\sum X_i$), $h=$ sabit çarpan. $\bar{X}_n$ yeterli.
\end{ornek}

\begin{ornek}[Çift parametreli üstel aile]\label{ex:cift_param_ustel}
Normal$(\mu,\sigma^2)$ her iki bilinmeyen: $\eta_1(\theta)=\mu/\sigma^2$, $\eta_2(\theta)=-1/(2\sigma^2)$. Doğal yeterli istatistik: $(\sum X_i, \sum X_i^2)$; tam rankli — tamamlayıcı.
\end{ornek}

\begin{teorem}[Yeterli istatistik ve bağımsızlık]\label{thm:yeterli_bagimsizlik}
$T$ yeterli ve $V$ herhangi bir istatistik. Eğer $\Prob_\theta(V\leq v\mid T=t)$ $\theta$'dan bağımsız ise, $V$'nin dağılımı yalnızca $T$'nin değerine bağlıdır. Özellikle $V$ bir pivot seçilirse çıkarım $T$ aracılığıyla yapılır.
\end{teorem}

% =====================================================================
\section{Bayes Çıkarımına Giriş}
\label{sec:bayes_giris}
% =====================================================================

\begin{kesif}
Frekansçı istatistikte $\theta$ sabit (bilinmeyen) bir değerdir; Bayesçi istatistikte ise $\theta$ da rasgele bir değişkendir ve ona dair önsel (prior) inancımızı bir olasılık dağılımı olarak ifade ederiz.
\end{kesif}

\begin{tanim}[Bayes teoremi — istatistik]\label{def:bayes_istatistik}
$\theta$'nın önsel dağılımı $\pi(\theta)$, veri olabilirliği $f(\mathbf{x}|\theta)$. \textbf{Sonsal dağılım} \emph{(posterior distribution)}:
\[
  \pi(\theta|\mathbf{x}) = \frac{f(\mathbf{x}|\theta)\pi(\theta)}{m(\mathbf{x})}, \quad m(\mathbf{x}) = \int f(\mathbf{x}|\theta)\pi(\theta)\,d\theta.
\]
$m(\mathbf{x})$ marjinal olabilirlik \emph{(marginal likelihood)} veya kanıt \emph{(evidence)}.
\end{tanim}

\begin{tanim}[Konjugat önsel]\label{def:konjugat_onsel}
$\pi(\theta)$ \textbf{konjugat önsel} \emph{(conjugate prior)}: sonsal $\pi(\theta|\mathbf{x})$ önsel ile aynı ailedendir. Hesap kolaylığı sağlar.
\begin{itemize}
  \item $\Binom(n,p)$ için $p\sim\Beta(\alpha,\beta)$: sonsal $\Beta(\alpha+\sum x_i, \beta+n-\sum x_i)$.
  \item $\Pois(\lambda)$ için $\lambda\sim\Gama(\alpha,\beta)$: sonsal $\Gama(\alpha+n\bar x, \beta+n)$.
  \item $\Normal(\mu,\sigma^2)$ ($\sigma^2$ bilinen) için $\mu\sim\Normal(\mu_0,\tau^2)$: sonsal $\Normal(\mu_1,\tau_1^2)$.
\end{itemize}
\end{tanim}

\begin{teorem}[Normal konjugat güncelleme]\label{thm:normal_konjugat}
$X_1,\ldots,X_n\iid\Normal(\mu,\sigma^2)$ ($\sigma^2$ bilinen); $\mu\sim\Normal(\mu_0,\tau^2)$. Sonsal:
\[
  \mu|\mathbf{X}\sim\Normal(\mu_1,\tau_1^2), \quad
  \mu_1 = \frac{\mu_0/\tau^2 + n\bar{X}/\sigma^2}{1/\tau^2+n/\sigma^2}, \quad
  \tau_1^2 = \frac{1}{1/\tau^2+n/\sigma^2}.
\]
\end{teorem}

\begin{proof}
$\pi(\mu|\mathbf{x})\propto\pi(\mu)f(\mathbf{x}|\mu)$. İkisi de Gaussian; çarpım da Gaussian (karede lineer). Üsteki $\exp[-(\mu-\mu_1)^2/(2\tau_1^2)+\text{sabit}]$ biçimine getirilince $\mu_1$ ve $\tau_1^2$ elde edilir.
\end{proof}

\begin{ornek}[Bayesçi ağırlıklı ortalama]
$\mu_1 = w_0\mu_0 + w_1\bar{X}$ ($w_0=(\sigma^2/n)/(\tau^2+\sigma^2/n)$, $w_1=\tau^2/(\tau^2+\sigma^2/n)$). Sonsal ortalama, önsel ortalama ile örneklem ortalamasının ağırlıklı ortalamasıdır. $n\to\infty$: $\mu_1\to\bar{X}$ (veri öneli baskılar); $\tau\to\infty$: $\mu_1\to\bar{X}$ (bilgisiz önsel).
\end{ornek}

\begin{teorem}[Bayes tahmincisi ve OKH]\label{thm:bayes_tahmin}
Karesel kayıp altında Bayes tahmincisi sonsal ortalama: $\hat\theta_{\text{Bayes}}=E[\theta|\mathbf{X}]$.

Mutlak değer kaybında Bayes tahmincisi sonsal medyan.
\end{teorem}

\begin{proof}
$E_{\theta|\mathbf{x}}[(\theta-d)^2]$ minimizasyonu: türev alıp sıfıra eşitle; $d^*=E[\theta|\mathbf{x}]$.
\end{proof}

\begin{kritiksorgu}
Frekansçı ve Bayesçi yaklaşım ne zaman aynı sonucu verir? Büyük örneklemde (Bernstein-von Mises teoremi) sonsal $\Normal(\hat\theta_{\mathrm{MLE}}, I(\hat\theta)^{-1}/n)$'ye yakınsar — Bayes ve frekansçı güven aralıkları örtüşür. Küçük örneklemde önsel seçim kritik önem taşır.
\end{kritiksorgu}

% =====================================================================
\section{Yeterli İstatistik ve Bayes}
\label{sec:yeterli_bayes}
% =====================================================================

\begin{teorem}[Yeterlilik ve Bayes]\label{thm:yeterlilik_bayes}
$T$ yeterli istatistik ise $\pi(\theta|\mathbf{x}) = \pi(\theta|T(\mathbf{x}))$. Yani Bayesçi güncelleme tamamen $T$ üzerinden gerçekleşir; orijinal verinin tüm $\theta$-ilgili bilgisi $T$'de toplanmıştır.
\end{teorem}

\begin{proof}
$\pi(\theta|\mathbf{x})\propto f(\mathbf{x}|\theta)\pi(\theta) = g(T(\mathbf{x}),\theta)h(\mathbf{x})\pi(\theta)$. $h(\mathbf{x})$ $\theta$'dan bağımsız; normalleştirmede iptal olur. Dolayısıyla $\pi(\theta|\mathbf{x})\propto g(T,\theta)\pi(\theta)=\pi(\theta|T)$.
\end{proof}

\begin{ornek}[Bayes ve yeterlilik — Bernoulli]
$X_i\iid\mathrm{Ber}(p)$; $T=\sum X_i$ yeterli. Sonsal $\pi(p|T=t)=\pi(p|\mathbf{x})$ — 100 gözlem yerine sadece $t=\sum x_i$ yeterli. Hiçbir ek bilgi kaybolmadı.
\end{ornek}

\begin{disiplinlerarasibaglan}
\textbf{Makine öğrenmesi:} Bayesçi sinir ağları (BNN) belirsizlik tahmini için sonsal dağılımı öğrenir. \textbf{Tıbbi teşhis:} Ön-test/son-test olasılığı (Bayes güncellemesi) tanısal testlerin yorumunda temel. \textbf{Sinyal işleme:} Kalman filtresi Bayesçi güncellemeler zinciridir (Gaussian konjugat).
\end{disiplinlerarasibaglan}

% =====================================================================
%  ALIŞTIRMALAR
% =====================================================================

\cozumlualistirmalar

\begin{alistirma}[Al.11.1]{A}{Örneklem Ortalaması ve Varyansı}
$X_1,\ldots,X_5$ i.i.d.\ $\Normal(3, 4)$ ($\sigma^2=4$).
(a) $\bar{X}_5$'in dağılımını yazın.
(b) $P(\bar{X}_5 > 4)$'ü hesaplayın.
(c) $(n-1)S^2/\sigma^2$'nin dağılımını yazın ve $E[S^2]$'yi bulun.
\end{alistirma}
\alcozum{%
(a) $\bar{X}_5 \sim \Normal(3, 4/5)$; $\sigma_{\bar{X}} = 2/\sqrt{5}$.\\
(b) $P(\bar{X}_5>4)=P(Z>(4-3)/(2/\sqrt{5}))=P(Z>\sqrt{5}/2)=P(Z>1.118)=1-\Phi(1.118)\approx0.132$.\\
(c) $4S^2/4 = S^2 \sim \chi^2_4/1$: $(n-1)S^2/\sigma^2 = 4S^2/4 \sim \chi^2_4$.
$E[S^2] = \sigma^2 = 4$.}

\begin{alistirma}[Al.11.2]{A}{Fisher-Neyman — Poisson}
$X_1,\ldots,X_n$ i.i.d.\ $\Pois(\lambda)$.
(a) Ortak PMF'yi yazın. (b) $T = \sum X_i$'nin yeterli istatistik olduğunu F-N teoremi ile gösterin.
(c) $T \sim \Pois(n\lambda)$ olduğunu belirtin.
\end{alistirma}
\alcozum{%
(a) $f(\mathbf{x}|\lambda) = \prod \frac{\lambda^{x_i}e^{-\lambda}}{x_i!} = \frac{\lambda^{\sum x_i}e^{-n\lambda}}{\prod x_i!}$.\\
(b) $g(T,\lambda) = \lambda^T e^{-n\lambda}$, $h(\mathbf{x}) = 1/\prod x_i!$ — ayrışım var.
F-N: $T = \sum X_i$ yeterli.\\
(c) Poisson yeniden üretkenliği: $T \sim \Pois(n\lambda)$.}

\begin{alistirma}[Al.11.3]{B}{OKH Hesabı}
$X_1,\ldots,X_n$ i.i.d.\ $\Normal(\mu,\sigma^2)$.
$\tilde{S}^2 = \frac{1}{n}\sum(X_i-\bar{X})^2$ (yanlı tahmin edici) için:
(a) $E[\tilde{S}^2]$ ve $\mathrm{Var}(\tilde{S}^2)$'yi bulun.
(b) $\mathrm{OKH}(\tilde{S}^2, \sigma^2)$'yi hesaplayın.
(c) $\mathrm{OKH}(\tilde{S}^2,\sigma^2)$ ile $\mathrm{OKH}(S^2,\sigma^2)$'yi karşılaştırın.
\end{alistirma}
\alcozum{%
(a) $\tilde{S}^2 = \frac{n-1}{n}S^2$; $E[\tilde{S}^2] = \frac{n-1}{n}\sigma^2$.
$\mathrm{Var}(\tilde{S}^2) = \left(\frac{n-1}{n}\right)^2\mathrm{Var}(S^2) = \left(\frac{n-1}{n}\right)^2 \frac{2\sigma^4}{n-1} = \frac{2(n-1)\sigma^4}{n^2}$.\\
(b) $\mathrm{Yan} = -\sigma^2/n$; $\mathrm{OKH}(\tilde{S}^2) = \frac{2(n-1)\sigma^4}{n^2}+\frac{\sigma^4}{n^2} = \frac{(2n-1)\sigma^4}{n^2}$.\\
(c) $\mathrm{OKH}(S^2) = \mathrm{Var}(S^2) = \frac{2\sigma^4}{n-1}$.
$n \geq 3$'te $\tilde{S}^2$ daha küçük OKH'ye sahip (küçük yanlılık kazancı).}

\begin{alistirma}[Al.11.4]{B}{Tamamlayıcılık — Uniform}
$X_1,\ldots,X_n$ i.i.d.\ $\Uniform(0,\theta)$.
(a) $T = X_{(n)} = \max X_i$'nin yeterli olduğunu F-N ile gösterin.
(b) $E_\theta[g(T)] = 0$ tüm $\theta > 0$'da koşulundan $g \equiv 0$ sonucunu türetin.
(c) $T$ tamamlayıcı mı?
\end{alistirma}
\alcozum{%
(a) $f(\mathbf{x}|\theta) = \theta^{-n}\mathbf{1}_{x_{(n)}\leq\theta}$;
$g(T,\theta) = \theta^{-n}\mathbf{1}_{T\leq\theta}$, $h=1$. F-N: $T$ yeterli.\\
(b) $f_{X_{(n)}}(t|\theta) = nt^{n-1}/\theta^n$ ($0<t<\theta$).
$E[g(T)] = \int_0^\theta g(t)\frac{nt^{n-1}}{\theta^n}dt = 0$.
$\frac{d}{d\theta}[\theta^n \cdot 0] = 0 = ng(\theta)\theta^{n-1}$, yani $g(\theta)=0$ tüm $\theta>0$.
(c) Evet, tamamlayıcı.}

% ---

\alistirmalar

\begin{alistirma}[Al.11.5]{A}{Örneklem $t$ Dağılımı}
$X_1,\ldots,X_{16}$ i.i.d.\ $\Normal(\mu, 9)$.
$T = \sqrt{16}(\bar{X}_{16}-\mu)/S$'nin dağılımını yazın.
$P(T > 2.131)$'i $t$ tablosundan söyleyin.
\end{alistirma}
\alcozum{%
$T \sim t_{15}$.
$t_{15,0.025} \approx 2.131$: $P(T > 2.131) = 0.025$.}

\begin{alistirma}[Al.11.6]{A}{Yeterli İstatistik — Üstel Aile}
$X_i \sim \Gama(\alpha, \beta)$ ($\alpha$ biliniyor, $\beta$ bilinmiyor).
(a) Ortak PDF'yi üstel aile formunda yazın. (b) Yeterli istatistiği bulun.
\end{alistirma}
\alcozum{%
(a) $f(\mathbf{x}|\beta) = \prod \frac{\beta^\alpha}{\Gamma(\alpha)}x_i^{\alpha-1}e^{-\beta x_i}
= \left(\frac{\beta^\alpha}{\Gamma(\alpha)}\right)^n \prod x_i^{\alpha-1} \cdot e^{-\beta\sum x_i}$.
(b) $g(T,\beta) = (\beta^\alpha/\Gamma(\alpha))^n e^{-\beta T}$, $h = \prod x_i^{\alpha-1}$.
$T = \sum X_i$ yeterli.}

\begin{alistirma}[Al.11.7]{B}{Bağımsızlık — $\bar{X}$ ve $S^2$}
$X_1, X_2, X_3$ i.i.d.\ $\Normal(\mu, \sigma^2)$.
$(Y_1, Y_2, Y_3) = (X_1+X_2+X_3, X_1-X_2, X_1+X_2-2X_3)/\sqrt{6}$ dönüşümünü yapın.
(a) $Y_i$'lerin $\Normal(?, ?)$ dağılımını bulun.
(b) $Y_1$ ile $(Y_2, Y_3)$'ün bağımsız olduğunu gösterin.
(c) $\bar{X}$ ve $S^2$'nin bağımsızlığını bu sonuçtan çıkarın.
\end{alistirma}
\alcozum{%
(a) $Y_1 = (X_1+X_2+X_3)/\sqrt{3}\cdot\sqrt{3}/\sqrt{6}\cdot\sqrt{6}$ — uygun normalize ile
$Y_1 = \sqrt{3}\bar{X}$: $\Normal(\sqrt{3}\mu, \sigma^2)$.
$Y_2 = (X_1-X_2)/\sqrt{2}$: $\Normal(0,\sigma^2)$. $Y_3 = (X_1+X_2-2X_3)/\sqrt{6}$: $\Normal(0,\sigma^2)$.\\
(b) $(Y_1,Y_2,Y_3)$ bağımsız normal: kovaryans matrisi köşegen (ortogonal dönüşüm).
$\mathrm{Cov}(Y_1,Y_2) = 0$; normal için $\Rightarrow$ bağımsız.\\
(c) $\bar{X} = Y_1/\sqrt{3}$; $S^2 = (Y_2^2+Y_3^2)/2\sigma^2 \cdot \sigma^2/1$
($(Y_2^2+Y_3^2)/\sigma^2 \sim \chi^2_2$). $Y_1 \perp (Y_2,Y_3)$ $\Rightarrow$ $\bar{X} \perp S^2$.}

\begin{alistirma}[Al.11.8]{B}{En Az Yeterlilik — Cauchy}
$X_1,\ldots,X_n$ i.i.d.\ $\mathrm{Cauchy}(\theta, 1)$.
$f(\mathbf{x}|\theta)/f(\mathbf{y}|\theta)$ oranını yazın.
Bu oranın $\theta$'dan bağımsız olması için hangi koşul gerekir?
En az yeterli istatistik nedir?
\end{alistirma}
\alcozum{%
$f(\mathbf{x}|\theta) = \prod \frac{1}{\pi(1+(x_i-\theta)^2)}$.
Oran: $\prod \frac{1+(y_i-\theta)^2}{1+(x_i-\theta)^2}$ — bu $\theta$'dan bağımsız olması için
$(x_1,\ldots,x_n)$ ve $(y_1,\ldots,y_n)$'nin aynı sıra istatistiklerine sahip olması gerekir
(çünkü Cauchy simetrik; terimlerin $\theta$-bağımlılığı birbirini ancak permütasyonla götürür).
En az yeterli istatistik: $(X_{(1)},\ldots,X_{(n)})$ — tam sıralı vektör.
Bu, Cauchy'de boyut azaltılmış bir özet olmamasını gösterir.}

\begin{alistirma}[Al.11.9]{A}{Sıra istatistiği dağılımı}
$X_1,\ldots,X_5\iid\Uniform[0,1]$. (a) $X_{(3)}$'ün yoğunluğunu yaz. (b) $E[X_{(3)}]$'ü bul. (c) $\Prob(X_{(5)}>0.9)$'u hesapla.
\end{alistirma}
\alcozum{%
(a) $f_{X_{(3)}}(x)=5!/(2!\cdot2!)x^2(1-x)^2\cdot1=30x^2(1-x)^2$ ($0<x<1$). (b) $E[X_{(3)}]=3/(5+1)=1/2$ (Uniform sıra istatistiği beklentisi $k/(n+1)$). (c) $F_{X_{(5)}}(x)=x^5$; $\Prob(X_{(5)}>0.9)=1-0.9^5=1-0.5905=0.4095$.%
}

\begin{alistirma}[Al.11.10]{B}{Ampirik dağılım ve KS testi}
$n=5$ gözlem: $1.2, 2.5, 0.8, 3.1, 1.7$. $H_0: F=\mathrm{Exp}(1)$ (yani $F_0(x)=1-e^{-x}$) için KS istatistiğini hesapla.
\end{alistirma}
\alcozum{%
Sıralı: $0.8,1.2,1.7,2.5,3.1$. $F_n(x_k)=k/5$. $F_0$: $F_0(0.8)=0.551$, $F_0(1.2)=0.699$, $F_0(1.7)=0.817$, $F_0(2.5)=0.918$, $F_0(3.1)=0.955$. $|F_n-F_0|$: $|0.2-0.551|=0.351$, $|0.4-0.699|=0.299$, $|0.6-0.817|=0.217$, $|0.8-0.918|=0.118$, $|1-0.955|=0.045$. $D_5=0.351$; $\sqrt{5}\cdot0.351=0.785<1.36$: $H_0$ reddedilemez.%
}

\begin{alistirma}[Al.11.11]{B}{Min için $n$ belirleme}
$X_i\iid\Exp(\lambda)$. $P(X_{(1)}>1)=0.9$ olmasını istiyorsunuz. $n$ kaç olmalı?
\end{alistirma}
\alcozum{%
$X_{(1)}\sim\Exp(n\lambda)$ (minimum üstel i.i.d. bağımsızsa). $P(X_{(1)}>1)=e^{-n\lambda}=0.9$; $n=-\log(0.9)/\lambda$. $\lambda=1$: $n=-\log(0.9)\approx0.105$; yani $n=1$ yetmez, tam tam değil — düz: $e^{-n}=0.9$; $n=0.105$ gerçek sayı. Yani tek gözlemde $\lambda=0.105$ ya da $n=1$, $\lambda=0.105$. Pratik: $\lambda=1$ ise $P(X_{(1)}>1)=e^{-n}=0.9$ için $n=\lceil-\log0.9\rceil=1$.%
}

% =====================================================================
\endinput
